% DFS 分支图：弧 1->2, 1->5, 2->3, 2->4, 5->4；从1出发，出邻居按编号递增。
% 发现/完成时刻：1:1/10, 2:2/7, 3:3/4, 4:5/6, 5:8/9；发现序 1,2,3,4,5；完成序 3,4,2,5,1。
\begin{tikzpicture}[font=\figurefont\fontsize{8}{10}\selectfont,>=Stealth,
  vtx/.style={circle,draw=FigureStroke,line width=.8pt,minimum size=6.5mm,inner sep=0pt,text=FigureInk,fill=white},
  tree/.style={draw=FigureBlue,line width=1.2pt,->},
  other/.style={draw=FigureMuted,line width=.8pt,dashed,->},
  tm/.style={font=\figurefont\fontsize{7.5}{9}\selectfont,text=FigureInk},
  cell/.style={draw=FigureStroke,line width=.6pt,minimum width=5.4mm,minimum height=4.4mm,inner sep=0pt,fill=FigureBlueFill},
  top/.style={cell,fill=FigureYellowFill},
  note/.style={font=\FigureNoteFont\fontsize{7.5}{9}\selectfont,text=FigureMuted}]
% ---------- (a) ----------
\begin{scope}
  \node[vtx] (v1) at (0,0) {1};
  \node[vtx] (v2) at (-0.9,-1.05) {2};
  \node[vtx] (v5) at (0.9,-1.05) {5};
  \node[vtx] (v3) at (-1.5,-2.15) {3};
  \node[vtx] (v4) at (0,-2.15) {4};
  \draw[tree] (v1)--(v2); \draw[tree] (v1)--(v5);
  \draw[tree] (v2)--(v3); \draw[tree] (v2)--(v4);
  \draw[other] (v5)--(v4);
  \node[tm,right=1pt of v1] {1/10};
  \node[tm,left=1pt of v2] {2/7};
  \node[tm,right=1pt of v5] {8/9};
  \node[tm,left=1pt of v3] {3/4};
  \node[tm,below=1pt of v4] {5/6};
  \node at (0,-3.55) {(a) 发现/完成时刻};
\end{scope}
% ---------- (b) ----------
\begin{scope}[xshift=2.15cm,yshift=-2.3cm]
  \def\w{0.64}
  \foreach \k/\ev/\st in {0/{发现\\1}/{1},1/{发现\\2}/{1,2},2/{发现\\3}/{1,2,3},3/{完成\\3}/{1,2},4/{发现\\4}/{1,2,4},
                          5/{完成\\4}/{1,2},6/{完成\\2}/{1},7/{发现\\5}/{1,5},8/{完成\\5}/{1},9/{完成\\1}/{}}{
    \node[tm,align=center] at (\k*\w,-0.62) {\ev};
    \foreach \x [count=\h from 0] in \st {
      \node[cell] at (\k*\w,0.05+\h*0.44) {\x};
    }
  }
  % 栈顶着色（逐列手工指定）
  \foreach \k/\h/\x in {0/0/1,1/1/2,2/2/3,3/1/2,4/2/4,5/1/2,6/0/1,7/1/5,8/0/1}{
    \node[top] at (\k*\w,0.05+\h*0.44) {\x};
  }
  \node[note] at (9*\w,0.05) {空};
  \draw[FigureStroke,line width=.5pt] (-0.35,-0.2)--(9*\w+0.35,-0.2);
  \node at (4.5*\w,-1.25) {(b) 每个事件后的栈（底在下，黄色为栈顶）};
  \node[note,anchor=west] at (-0.35,1.55) {完成序：3, 4, 2, 5, 1};
  \node[note,anchor=west] at (-0.35,1.95) {发现序：1, 2, 3, 4, 5};
\end{scope}
\end{tikzpicture}
